%Acknowledgements are the author's statement of gratitude to and
%recognition of the people and institutions who helped the author's
%research and writing.

\begin{center}
\vspace*{52pt}
{\normalfont \textbf{ACKNOWLEDGEMENTS}}
\end{center}

This document represents the culmination of years of work, persistence, and a remarkable ability to continue moving forward despite frequent uncertainty about what, exactly, was going on.

If you are reading this as you prepare your own dissertation, it is worth noting that very few people feel entirely ready for this moment. The process is rarely linear, clarity is often delayed, and confidence tends to arrive well after it would have been most convenient. That is normal. Progress, more often than not, looks like persistence disguised as confusion.

There will be times when the work feels too large, the questions too complex, and the path forward unclear. In those moments, it may help to remember that the goal is not perfection, but completion—and that meaningful work is almost always the result of sustained effort rather than sudden insight.

You will likely revise more than you expect, doubt more than you anticipated, and learn far more than is immediately visible on the page. Somewhere along the way, the pieces will begin to cohere. And even if it does not feel like it, you are, in fact, becoming an expert.

So keep going. Trust the process just enough to continue. Ask good questions, accept imperfect answers, and write anyway.

And when you finally reach the end, take a moment to appreciate what you have done—because finishing a dissertation is, objectively, a slightly unreasonable thing to accomplish.

\clearpage
